\subsection{Savarankiškai sukurti ELF testai}
\label{sec:savarankiski-elf-testai}

Be oficialių RISC-V testų, projekte yra savarankiškai sukurti ELF testai. Jie skirti tikrinti privilegijuoto vykdymo dalis, kurios tiesiogiai reikalingos Linux paleidimui: \texttt{mtvec}/\texttt{stvec} vektorius, \texttt{mret}/\texttt{sret} instrukcijas, \texttt{mstatus}/\texttt{sstatus} laukus, išimčių delegavimą, \texttt{ecall}, neleistinas instrukcijas ir blogas atminties prieigas. Testas surenkamas kaip savarankiškas \texttt{rv64g} ELF failas.

Testai apima šiuos atvejus:

\begin{itemize}
    \item \textbf{neleistiną instrukciją} -- vykdoma \texttt{.word 0xffffffff} instrukcija, kuri turi sukelti illegal instruction išimtį;
    \item \textbf{sisteminius kvietimus} -- \texttt{ecall} tikrinamas M, S ir U režimuose;
    \item \textbf{privilegijų keitimą} -- \texttt{mret} naudojamas perėjimui iš M režimo į S arba U režimą, o \texttt{sret} -- perėjimui iš S režimo į U režimą;
    \item \textbf{išimčių delegavimą} -- \texttt{medeleg} registre įjungiama dalies išimčių delegacija S režimui;
    \item \textbf{blogas atminties prieigas} -- atliekamas skaitymas ir rašymas į neteisingą adresą, kad būtų patikrinta magistralės klaidų ir RISC-V išimčių sąsaja.
\end{itemize}

Testo sėkmė ir nesėkmė perduodama per \texttt{tohost} mechanizmą, apibrėžtą oficialiuose RISC-V testuose. Papildomai testai rašo diagnostinę informaciją į UART adresą \texttt{0x10000000}, todėl galima matyti, kuriame žingsnyje įvyko klaida ir kokios CSR reikšmės buvo tuo metu.

Šie savarankiški testai buvo naudingi todėl, kad jie tikrina ne vien instrukcijų rezultatą, bet ir visą trap kelią: privilegijos režimo pakeitimą, \texttt{mcause}/\texttt{scause} nustatymą, \texttt{mepc}/\texttt{sepc} grįžimo adresus ir trap handlerio pakartotinį įėjimą į C testų ciklą. Būtent šios dalys yra kritinės prieš pereinant prie Linux paleidimo.
